\section{Summary}
\begin{frame}{}
    \LARGE DQN Family \& SARSA: \textbf{Summary}
\end{frame}

\begin{frame}{When to Use Which Algorithm?}
    \begin{table}[]
        \centering
        \renewcommand{\arraystretch}{2.5}
        \begin{tabular}{ll}
            \hline
            \textbf{Scenario} & \textbf{Recommended Algorithm} \\ \hline
            Environment is stochastic & SARSA (conservative) \\ 
            Maximizing reward is the priority & Q-Learning or DQN \\ 
            Large state space (e.g., images) & DQN \\ 
            Limited computational resources & SARSA (simpler model) \\ \hline
        \end{tabular}
    \end{table}
\end{frame}

\begin{frame}{Future Directions}
    \small
    DQN works well but has known weaknesses. Each extension below addresses one specific problem. Together they form the DQN family.
    \begin{itemize}
        \setlength{\itemsep}{0.4em}
        \item \textbf{Dueling DQN}: Splits $Q$ into $V(s)$ and $A(s,a)$---learns which states are valuable without evaluating every action.
        \item \textbf{Prioritized Replay}: Samples transitions proportional to TD error instead of uniformly---prioritises what the network got wrong.
        \item \textbf{Distributional RL}: Learns the full return distribution $Z(s,a)$, not just the expected value---richer signal, faster learning.
        \item \textbf{Noisy Nets}: Replaces $\varepsilon$-greedy with learned noise in the weights---the network decides how much to explore.
        \item \textbf{Rainbow}: Combines all of the above. Represents the practical ceiling of value-based methods.
    \end{itemize}
\end{frame}

\begin{frame}[allowframebreaks]{Summary}
    \begin{itemize}
        \setlength{\itemsep}{1em}
        \item SARSA is an on-policy alternative to Q-learning, using the action the agent actually takes rather than $\max_{a'} Q(s',a')$.
        \item The on-policy vs.\ off-policy distinction is one of the most important axes in RL algorithm design.
        \item Deep Q-Learning uses a neural network as a function approximator for Q-values, enabling RL on high-dimensional inputs.
        \item Experience replay breaks correlation between consecutive samples; the target network stabilizes the moving target.
        \item Double DQN decouples action selection from evaluation to reduce overestimation bias.
        \item The DQN family extensions (Dueling, Prioritized Replay, Distributional, Noisy Nets, Rainbow) each address a specific weakness of vanilla DQN.
    \end{itemize}
\end{frame}